% \section{Lessons for Scientific Data Infrastructure}

% CARP demonstrates that adaptive range partitioning, driven by periodic reconstruction of global
% histograms, can effectively range partition scientific data in situ. The resulting data layout is
% partially ordered, requiring some residual compaction at query time. Our evaluation demonstrates
% that this is the right tradeoff for scientific data: differences in query performance are
% imperceptible for analytical usecases, while ingestion performance is maximized.


% 1. Scientific data pipelines ingest synchronously from a massive stream of parallel writers,
% terminating in a parallel filesystem. Parallel filesystems and burst buffers are designed to
% absorb this synchronous and continuous burst. While database-style approaches like range
% partitioning remain deeply relevant, they need to be rearchitected for this environment. Partially
% organized data lyouts, persisted as append-only logs to a filesystem, are a natural fit for this
% environment. The goal is to uplevel database techniques into this pipeline, rather than ingesting
% into a database itself.

% 2. Similarity with OLAP. our storage backend implements a custom format that derives from
% DeltaFS's flattened LSM tree: a set of disjoint SSTables. It implements contiguous key and value
% blocks for more efficient sequential reads. Statistics over overlapping ranges are maintained in a
% catalog, and used to prune irrelevant partitions at query time. This is essentially identical to
% Parquet and Iceberg, and our first brush with analytics bottlenecks and columnar remedies that
% shows up over the rest of this thesis. Combined with the previous point, it further underscores
% the natural fit of this paradigm for scientific data pipelines. A natural concern is
% interoperability/composability between CARP's partitioning and additional indexing strategies,
% which are automatically remedied by planning-based approaches. We discuss this more in
% (thesis-conclusions).

% 3. Limitations of a fixed-function loop. CARP's loop can materialize only one view using one
% reduction technique. Additional views are desirable not only for other workload types, but also
% for better range partitioning. For example, sketch-based quintiles may provide better distribution
% fidelity over histograms. Renegotiation policies may need to be more adaptive: CARP relies on more
% frequent renegotiation not being counterproductive, but this may not hold.

\section{Lessons for Scientific Data Infrastructure}
\label{sec:carp:lessons}

CARP demonstrates that adaptive range partitioning, driven by periodic reconstruction of global
histograms, can effectively range-partition scientific data in situ.
The resulting layout is
partially ordered, so range queries may require some residual compaction at query time.
Our
evaluation shows that this is the right tradeoff for scientific data: differences in query
performance are
imperceptible for analytical use cases, while ingestion performance is maximized.
% binding constraint, while differences in query performance are not material for analytical use
% cases. 
Two lessons recur across the AMR study discussed next and inform ORCA’s design: the
applicability of OLAP techniques to scientific data management, and the limits of
fixed-function feedback loops.

\paragraph{Applicability of OLAP techniques} While hash and range-partitioning are standard
database
techniques, their application to scientific data management is shaped by the environment.
These pipelines terminate in parallel filesystems~\cite{Braam2002:Lustre,Oeste2025:StorageGaps} designed to absorb
massive,
synchronized ingestion from tens of thousands of writers.
As a result, these techniques must be embedded into in-situ pipelines
that reorganize data in flight to parallel filesystems.
% While techniques remain deeply relevant, they need to be rearchitected for in-situ pipelines that
% reorganize data in flight to a parallel filesystem.

Beyond techniques, OLAP formats also remain relevant for scientific data, at least insofar as
data fits the relational model.
CARP evolves a custom storage format derived from DeltaFS's~\cite{Zheng2018:DeltaFS}
flattened LSM tree: disjoint SSTables with contiguous key and value blocks for efficient sequential
reads,
with per-SSTable statistics maintained in a catalog for query-time pruning.
This is essentially identical to Parquet~\cite{:ApacheParquet} and open table formats like
Iceberg~\cite{:Iceberg}.
Queryability over other attributes is a natural concern automatically remedied by such
formats
and planning-based query engines.
We revisit this topic in \S\ref{sec:concl:extensions}.

\paragraph{Limitations of fixed-function loops} CARP’s feedback loop is fixed-function: it can
materialize one view via one reduction type, triggered by a static policy.
Additional views
may be desirable not only for other workload types but for better range partitioning itself.
Sketch-based quantile estimates~\cite{Karni2016:KLLSketch} may provide better distribution fidelity
than histograms.
Decisions such as tuning renegotiation intervals and pivot counts would
benefit from runtime visibility into how well the partitioning is performing---visibility that
CARP’s own loop cannot provide.

% Database and analytics patterns recur throughout scientific data management, but the environment
% constrains how they apply. Workloads are diverse and configurations are transient, so no single
% database architecture fits as the default endpoint. Ingestion to a database is not viable. The
% pipeline itself must be the site of data management. Scientific data pipelines terminate in a
% parallel filesystem or burst buffer that is designed to absorb massive, synchronized ingestion
% from thousands of writers. Workloads are diverse and configurations are transient, so no single
% database architecture fits as the default endpoint. Ingestion to a database is not viable. The
% pipeline itself must be the site of data management.

% What translates are techniques and formats. CARP demonstrates range partitioning adapted into an
% ingestion pipeline, operating over data in flight rather than reorganizing it at rest. The storage
% format converges in the same direction. Our backend derives from DeltaFS’s flattened LSM tree:
% disjoint SSTables with contiguous key and value blocks for efficient sequential reads, and
% per-range statistics maintained in a catalog for query-time pruning. This is essentially Parquet
% with an Iceberg-style catalog. This convergence between scientific data management and analytical
% storage recurs throughout the thesis and is revisited in the conclusions.

% Renegotiation policies may need to be more adaptive. Other views may enable adaptive range
% partitioning policies, or guidance to operators choosing renegotiation intervals and the fidelity
% of statistics. CARP assumes that more frequent renegotiation is not counterproductive, but this
% may not hold in all settings.
% Tuning these decisions would
